\documentclass[a4paper,landscape]{article}
\usepackage[a4paper,margin=0.4in]{geometry} % Adjust margins to maximize space
\usepackage{multicol}
\usepackage{titlesec}
\usepackage{enumitem}
\usepackage{amsmath, amssymb}
\usepackage{hyperref}
\usepackage{setspace}
\usepackage{graphicx}

% Small text settings
\renewcommand{\baselinestretch}{0.75}
\setlength{\columnsep}{1cm}

% Title formatting
\titleformat{\section}{\large\bfseries}{}{0em}{}[\titlerule]
\titleformat{\subsection}{\normalsize\bfseries}{}{0em}{}
\titleformat{\subsubsection}[runin]{\bfseries}{}{0em}{}[:]

\begin{document}

\raggedright
\footnotesize
\begin{multicols*}{5}

% Title
\section*{MidTerm2 MGLI Cheatsheet Doruk}

% Part 1: Spanning Trees
\section*{Part 1: Spanning Trees}

\subsection*{Core Definitions}
\begin{itemize}[leftmargin=*, itemsep=0pt]
    \item \textbf{Tree:} Connected, acyclic graph
    \begin{itemize}[leftmargin=*, itemsep=0pt]
        \item Unique path between nodes
        \item $n$ nodes → $n-1$ edges
        \item Connected + acyclic
    \end{itemize}
    \item \textbf{Leaf:} Node with degree 1
    \item \textbf{Spanning Tree:} Tree connecting all nodes with min edges
    \item \textbf{Edge-Weighted:} Edges have numeric weights
\end{itemize}

\subsection*{Tree Properties}
\begin{itemize}[leftmargin=*, itemsep=0pt]
    \item \textbf{Unique Path:} Only one path between any nodes
    \item \textbf{Edge Count:} $n$ nodes → $n-1$ edges
    \item \textbf{Minimal:} Remove edge → disconnected
    \item \textbf{Maximal:} Add edge → creates cycle
\end{itemize}

\subsection*{MST (Minimum Spanning Tree)}
\begin{itemize}[leftmargin=*, itemsep=0pt]
    \item \textbf{Definition:} Spanning tree with min total weight
    \item \textbf{Key Points:} Non-unique, all vertices, min weight sum
\end{itemize}

\subsection*{Greedy MST Algorithms}
\textbf{Kruskal's Algorithm Steps:}
\begin{itemize}[leftmargin=*, itemsep=0pt]
    \item Sort edges by weight ascending
    \item Process each edge in order
    \item Add edge if it doesn't create cycle
    \item Stop when $n-1$ edges added
\end{itemize}

\textbf{Prim's Algorithm Steps:}
\begin{itemize}[leftmargin=*, itemsep=0pt]
    \item Start from any vertex
    \item Add minimum weight edge to tree
    \item Only consider edges connecting to unvisited vertices
    \item Continue until all vertices included
\end{itemize}

\subsection*{Graph Components}
\begin{itemize}[leftmargin=*, itemsep=0pt]
    \item \textbf{Connected Component:} Maximal connected subgraph
    \item \textbf{Edge Effects:}
    \begin{itemize}[leftmargin=*, itemsep=0pt]
        \item Removing edge from tree: Creates two components
        \item Adding edge between components: Creates new spanning tree
    \end{itemize}
\end{itemize}

\subsection*{Advanced Concepts}
\begin{itemize}[leftmargin=*, itemsep=0pt]
    \item \textbf{Weighted Graph:} $w: E \rightarrow \mathbb{R}$ assigns weights
    \item \textbf{Tree Cost:} Sum of all edge weights in tree
    \item \textbf{Steiner Trees:} Minimal trees including specified vertex subset
    \item \textbf{Depth-First Search:} Method to explore tree structure
\end{itemize}
\subsection*{Key Properties}
\begin{itemize}[leftmargin=*, itemsep=0pt]
    \item \textbf{Number of Possible Trees (Cayley's Formula):} In a graph where every vertex connects to every other vertex by edges, there are $n^{n-2}$ different possible tree arrangements, where $n$ is the number of vertices. For example, with 4 vertices, there are $4^2 = 16$ different possible trees you can make.
    \item \textbf{Tree Encoding (Prüfer Code):} A Prüfer Code is a clever way to represent a tree using a sequence of numbers. For a tree with $n$ vertices, this code always has length $n-2$. For instance, a tree with 5 vertices can be encoded in a sequence of 3 numbers, making it easier to store and work with tree structures.
    \item \textbf{Prüfer Code Example:}
    \vspace*{5cm}
\end{itemize}

% Add additional spanning tree content here as needed

% Part 2: Propositional Logic
\section*{Part 2: Propositional Logic}

\subsection*{Core Concepts}
\begin{itemize}[leftmargin=*, itemsep=0pt]
    \item \textbf{Proposition:} Declarative sentence, true or false.
    \item \textbf{Atomic Propositions:} Basic, indivisible.
\end{itemize}

\subsection*{Logical Operators}
\begin{itemize}[leftmargin=*, itemsep=0pt]
    \item \textbf{Conjunction $(\wedge)$:} True if both operands true.
    \item \textbf{Disjunction $(\vee)$:} True if at least one operand true.
    \item \textbf{Negation $(\neg)$:} Inverts truth value.
\end{itemize}

\subsection*{Normal Forms}
\begin{itemize}[leftmargin=*, itemsep=0pt]
    \item \textbf{CNF:} Conjunction of disjunctions.
    \item \textbf{DNF:} Disjunction of conjunctions.
\end{itemize}

\subsection*{Normal Forms Construction}
\begin{itemize}[leftmargin=*, itemsep=0pt]
    \item \textbf{CNF Steps:}
    \begin{itemize}[leftmargin=*, itemsep=0pt]
        \item Identify false rows in truth table
        \item Create clause for each false row
        \item Combine clauses with AND
    \end{itemize}
    \item \textbf{Example CNF:} For formula $(A \rightarrow B) \wedge \neg C$:
    \begin{itemize}[leftmargin=*, itemsep=0pt]
        \item False rows: $(A=1,B=0,C=0), (A=1,B=0,C=1), (A=0,B=0,C=1), (A=0,B=1,C=1), (A=1,B=1,C=1)$
        \item Clauses: $(\neg A \vee B)$ for $A \rightarrow B$ and $(\neg C)$ for $\neg C$
        \item CNF: $(\neg A \vee B) \wedge (\neg C)$
    \end{itemize}
    \item \textbf{DNF Steps:}
    \begin{itemize}[leftmargin=*, itemsep=0pt]
        \item Identify true rows in truth table
        \item Create minterm for each true row
        \item Combine minterms with OR
    \end{itemize}
    \item \textbf{Example DNF:} For formula $A \vee B$:
    \begin{itemize}[leftmargin=*, itemsep=0pt]
        \item True rows: $(A=0,B=1), (A=1,B=0), (A=1,B=1)$
        \item Minterms: $(\neg A \wedge B), (A \wedge \neg B), (A \wedge B)$
        \item DNF: $(\neg A \wedge B) \vee (A \wedge \neg B) \vee (A \wedge B)$
    \end{itemize}
\end{itemize}

\subsection*{Truth Tables}
\begin{center}
\begin{tabular}{cc|c}
A & B & A$\wedge$B \\
\hline
0 & 0 & 0 \\
0 & 1 & 0 \\
1 & 0 & 0 \\
1 & 1 & 1
\end{tabular}
\quad
\begin{tabular}{cc|c}
A & B & A$\vee$B \\
\hline
0 & 0 & 0 \\
0 & 1 & 1 \\
1 & 0 & 1 \\
1 & 1 & 1
\end{tabular}
\end{center}

\subsection*{Statement Properties}
\begin{itemize}[leftmargin=*, itemsep=0pt]
    \item \textbf{Valid Statements:}
    \begin{itemize}[leftmargin=*, itemsep=0pt]
        \item Must have definite truth value
        \item Must be declarative
        \item Cannot be self-contradictory
    \end{itemize}
    \item \textbf{Invalid Examples:}
    \begin{itemize}[leftmargin=*, itemsep=0pt]
        \item Questions
        \item Commands
        \item Self-referential statements
    \end{itemize}
\end{itemize}

\subsection*{Logical Formula Structure}
\begin{itemize}[leftmargin=*, itemsep=0pt]
    \item \textbf{Syntax:} Rules for valid formula construction
    \item \textbf{Semantics:} Meaning and evaluation rules
    \item \textbf{Variable Assignment:} Maps variables to truth values
    \item \textbf{Evaluation:} Computes formula's truth value
\end{itemize}

\subsection*{Formula Properties}
\begin{itemize}[leftmargin=*, itemsep=0pt]
    \item \textbf{Satisfiable:} True for some assignment
    \item \textbf{Tautology:} True for all assignments
    \item \textbf{Contradiction:} False for all assignments
    \item \textbf{SAT Problem:} Finding satisfying assignment
\end{itemize}

\subsection*{Important Laws}
\begin{itemize}[leftmargin=*, itemsep=0pt]
    \item \textbf{De Morgan:} 
    \begin{itemize}[leftmargin=*, itemsep=0pt]
        \item $\neg(A\wedge B) \equiv \neg A\vee\neg B$
        \item $\neg(A\vee B) \equiv \neg A\wedge\neg B$
    \end{itemize}
    \item \textbf{Distributive:}
    \begin{itemize}[leftmargin=*, itemsep=0pt]
        \item $A\wedge(B\vee C) \equiv (A\wedge B)\vee(A\wedge C)$
        \item $A\vee(B\wedge C) \equiv (A\vee B)\wedge(A\vee C)$
    \end{itemize}
\end{itemize}

\subsection*{Additional Laws}
\begin{itemize}[leftmargin=*, itemsep=0pt]
    \item \textbf{Idempotence:}
    \begin{itemize}[leftmargin=*, itemsep=0pt]
        \item $A \wedge A \equiv A$
        \item $A \vee A \equiv A$
    \end{itemize}
    \item \textbf{Absorption:}
    \begin{itemize}[leftmargin=*, itemsep=0pt]
        \item $A \wedge (A \vee B) \equiv A$
        \item $A \vee (A \wedge B) \equiv A$
    \end{itemize}
    \item \textbf{Double Negation:}
        \item $\neg\neg A \equiv A$
    \item \textbf{Identity:}
    \begin{itemize}[leftmargin=*, itemsep=0pt]
        \item $A \wedge 1 \equiv A$
        \item $A \vee 0 \equiv A$
        \item $A \wedge 0 \equiv 0$
        \item $A \vee 1 \equiv 1$
    \end{itemize}
\end{itemize}

% Part 3: Predicates and Quantifiers
\section*{Part 3: Predicates and Quantifiers}

\subsection*{Implication and Equivalence}
\begin{itemize}[leftmargin=*, itemsep=0pt]
    \item \textbf{Implication $(\Rightarrow)$:} If $A$ is true, then $B$ is true.
    \begin{itemize}[leftmargin=*, itemsep=0pt]
        \item Notation: $A \Rightarrow B$
        \item Truth Table:
        \begin{center}
        \begin{tabular}{cc|c}
        A & B & $A \Rightarrow B$ \\
        \hline
        0 & 0 & 1 \\
        0 & 1 & 1 \\
        1 & 0 & 0 \\
        1 & 1 & 1 \\
        \end{tabular}
        \end{center}
        \item Example: "If it rains, the ground is wet."
    \end{itemize}
    \item \textbf{Equivalence $(\Leftrightarrow)$:} $A$ is true if and only if $B$ is true.
    \begin{itemize}[leftmargin=*, itemsep=0pt]
        \item Notation: $A \Leftrightarrow B$
        \item Truth Table:
        \begin{center}
        \begin{tabular}{cc|c}
        A & B & $A \Leftrightarrow B$ \\
        \hline
        0 & 0 & 1 \\
        0 & 1 & 0 \\
        1 & 0 & 0 \\
        1 & 1 & 1 \\
        \end{tabular}
        \end{center}
        \item Example: "A number is even if and only if it is divisible by 2."
    \end{itemize}
\end{itemize}

\subsection*{Contraposition}
\begin{itemize}[leftmargin=*, itemsep=0pt]
    \item \textbf{Principle:} $A \Rightarrow B$ is equivalent to $\neg B \Rightarrow \neg A$
    \item Example: "If I drink coffee, I am awake." Contraposition: "If I am not awake, I did not drink coffee."
\end{itemize}

\subsection*{Predicates}
\begin{itemize}[leftmargin=*, itemsep=0pt]
    \item \textbf{Definition:} Statement with variables, becomes true/false when variables are assigned values.
    \item \textbf{Examples:} $P(x):$ "$x$ is a prime number."
\end{itemize}

\subsection*{Quantifiers}
\begin{itemize}[leftmargin=*, itemsep=0pt]
    \item \textbf{Universal Quantifier $(\forall):$} True for all elements in domain.
    \begin{itemize}[leftmargin=*, itemsep=0pt]
        \item Notation: $\forall x \in M, P(x)$
        \item Example: "All humans are mortal."
    \end{itemize}
    \item \textbf{Existential Quantifier $(\exists):$} True for at least one element in domain.
    \begin{itemize}[leftmargin=*, itemsep=0pt]
        \item Notation: $\exists x \in M, P(x)$
        \item Example: "There exists a prime number greater than 2."
    \end{itemize}
\end{itemize}

\subsection*{Quantifier Rules}
\begin{itemize}[leftmargin=*, itemsep=0pt]
    \item \textbf{Negation:}
    \begin{itemize}[leftmargin=*, itemsep=0pt]
        \item $\neg(\forall x P(x)) \equiv \exists x \neg P(x)$
        \item $\neg(\exists x P(x)) \equiv \forall x \neg P(x)$
    \end{itemize}
    \item \textbf{Distribution:}
    \begin{itemize}[leftmargin=*, itemsep=0pt]
        \item $\forall x(P(x) \wedge Q(x)) \equiv (\forall xP(x)) \wedge (\forall xQ(x))$
        \item $\exists x(P(x) \vee Q(x)) \equiv (\exists xP(x)) \vee (\exists xQ(x))$
    \end{itemize}
\end{itemize}

\subsection*{Example}
\begin{itemize}[leftmargin=*, itemsep=0pt]
    \item "All apples are red" can be expressed as $\forall x \in \text{Apples}, \text{Red}(x)$.
\end{itemize}


\section*{Part 1: Quantifier Rules \& Properties}

\subsection*{Quantifier Order}
\begin{itemize}[leftmargin=*, itemsep=0pt]
    \item Different quantifiers cannot be reordered:
    \begin{itemize}[leftmargin=*, itemsep=0pt]
        \item $\forall x \exists y P(x,y) \not\equiv \exists y \forall x P(x,y)$
        \item Example: "For every person, there exists a film they know" $\not\equiv$ "There exists a film that every person knows"
    \end{itemize}
    \item Same quantifiers can be reordered:
    \begin{itemize}[leftmargin=*, itemsep=0pt]
        \item $\forall x \forall y P(x,y) \equiv \forall y \forall x P(x,y)$
        \item $\exists x \exists y P(x,y) \equiv \exists y \exists x P(x,y)$
    \end{itemize}
\end{itemize}

\subsection*{Distributive Laws for Quantifiers}
\textbf{Existential Quantifier $(\exists)$:}
\begin{itemize}[leftmargin=*, itemsep=0pt]
    \item Distributes over OR:
        \item $\exists x(P(x) \vee Q(x)) \equiv (\exists x P(x)) \vee (\exists x Q(x))$
    \item Does NOT distribute over AND:
        \item $\exists x(P(x) \wedge Q(x)) \not\equiv (\exists x P(x)) \wedge (\exists x Q(x))$
    \item But: $\exists x(P(x) \wedge Q(x)) \Rightarrow (\exists x P(x)) \wedge (\exists x Q(x))$
\end{itemize}

\textbf{Universal Quantifier $(\forall)$:}
\begin{itemize}[leftmargin=*, itemsep=0pt]
    \item Distributes over AND:
        \item $\forall x(P(x) \wedge Q(x)) \equiv (\forall x P(x)) \wedge (\forall x Q(x))$
    \item Does NOT distribute over OR:
        \item $\forall x(P(x) \vee Q(x)) \not\equiv (\forall x P(x)) \vee (\forall x Q(x))$
    \item But: $(\forall x P(x)) \vee (\forall x Q(x)) \Rightarrow \forall x(P(x) \vee Q(x))$
\end{itemize}

\section*{Part 2: Summation \& Product Notation}

\subsection*{Summation Rules}
\begin{itemize}[leftmargin=*, itemsep=0pt]
    \item \textbf{Constant Rule:} $\sum_{i=m}^n c = (n-m+1)c$
    \item \textbf{Scalar Rule:} $\sum_{i=m}^n ca_i = c\sum_{i=m}^n a_i$
    \item \textbf{Split Sum:} $\sum_{i=m}^n a_i = \sum_{i=m}^k a_i + \sum_{i=k+1}^n a_i$
    \item \textbf{Sum Addition:} $\sum_{i=m}^n (a_i + b_i) = \sum_{i=m}^n a_i + \sum_{i=m}^n b_i$
\end{itemize}

\subsection*{Product Rules}
\begin{itemize}[leftmargin=*, itemsep=0pt]
    \item \textbf{Constant Rule:} $\prod_{i=m}^n c = c^{n-m+1}$
    \item \textbf{Power Rule:} $\prod_{i=m}^n c^{a_i} = c^{\sum_{i=m}^n a_i}$
    \item \textbf{Split Product:} $\prod_{i=m}^n a_i = (\prod_{i=m}^k a_i)(\prod_{i=k+1}^n a_i)$
    \item \textbf{Product Multiplication:} $\prod_{i=m}^n (a_i \cdot b_i) = (\prod_{i=m}^n a_i)(\prod_{i=m}^n b_i)$
\end{itemize}

\section*{Part 3: Mathematical Induction}

\subsection*{Principle of Mathematical Induction}
\begin{itemize}[leftmargin=*, itemsep=0pt]
    \item \textbf{Base Case:} Prove $P(n_0)$ is true
    \item \textbf{Inductive Step:} Prove $P(k) \Rightarrow P(k+1)$ for $k \geq n_0$
    \item Base case must be proven - cannot be skipped
    \item Base case can't start after $n_0$
\end{itemize}

\subsection*{Common Induction Applications}
\begin{itemize}[leftmargin=*, itemsep=0pt]
    \item \textbf{Summation Formulas:}
    \begin{itemize}[leftmargin=*, itemsep=0pt]
        \item Arithmetic series: $\sum_{i=1}^n i = \frac{n(n+1)}{2}$
        \item Geometric series: $\sum_{i=0}^n 2^i = 2^{n+1} - 1$
    \end{itemize}
    \item \textbf{Inequalities:}
    \begin{itemize}[leftmargin=*, itemsep=0pt]
        \item Bernoulli: $(1+x)^n \geq 1 + nx$ for $x \geq -1$
    \end{itemize}
    \item \textbf{Divisibility:}
    \begin{itemize}[leftmargin=*, itemsep=0pt]
        \item Example: $8^n - 1$ is divisible by 7
    \end{itemize}
\end{itemize}

\subsection*{Induction Tips}
\begin{itemize}[leftmargin=*, itemsep=0pt]
    \item Identify pattern before attempting proof
    \item Write out first few cases
    \item In inductive step, assume $P(k)$ before proving $P(k+1)$
    \item Look for ways to use inductive hypothesis
\end{itemize}

\section*{Part 4: Common Examples \& Applications}

\subsection*{Key Proofs by Induction}
\begin{itemize}[leftmargin=*, itemsep=0pt]
    \item \textbf{Sum of First n Numbers:}
    \begin{itemize}[leftmargin=*, itemsep=0pt]
        \item $1 + 2 + ... + n = \frac{n(n+1)}{2}$
        \item Base: $n=1$: $1 = \frac{1(2)}{2}$
        \item Step: Add $(k+1)$ to both sides
    \end{itemize}
    \item \textbf{Geometric Series:}
    \begin{itemize}[leftmargin=*, itemsep=0pt]
        \item $\sum_{i=0}^n 2^i = 2^{n+1} - 1$
        \item Base: $n=0$: $2^0 = 1 = 2^1 - 1$
        \item Step: Add $2^{k+1}$ to both sides
    \end{itemize}
\end{itemize}

\section*{Part 4: Recursion \& Sequences}

\subsection*{Core Concepts}
\begin{itemize}[leftmargin=*, itemsep=0pt]
    \item \textbf{Recursion:} From Latin "recurrere" (to run back)
    \item Structure defined in terms of itself
    \item Self-similar patterns in nature and mathematics
    \item Examples: Fractals, Recursive sequences
\end{itemize}

\subsection*{Sequence Representations}
\begin{itemize}[leftmargin=*, itemsep=0pt]
    \item \textbf{Formal Definition:} Function $f: \mathbb{N} \rightarrow \mathbb{R}$
        \begin{itemize}[leftmargin=*, itemsep=0pt]
            \item Example: $f(n) = 2n + 1$ maps natural numbers to odd numbers
        \end{itemize}
    \item \textbf{Enumerated Form:} List all terms $f(1), f(2), f(3), ...$
        \begin{itemize}[leftmargin=*, itemsep=0pt]
            \item Example: Powers of 2: 2, 4, 8, 16, 32, ...
            \item Example: Perfect squares: 1, 4, 9, 16, 25, ...
        \end{itemize}
    \item \textbf{Explicit Form:} Direct formula for $f(n)$
        \begin{itemize}[leftmargin=*, itemsep=0pt]
            \item Example: $f(n) = n^2$ for perfect squares
            \item Example: $f(n) = 2^n$ for powers of 2
        \end{itemize}
    \item \textbf{Recursive Form:} Define in terms of previous terms
        \begin{itemize}[leftmargin=*, itemsep=0pt]
            \item Example: $f(1)=2$, $f(n)=2f(n-1)$ for powers of 2
            \item Example: $f(1)=1$, $f(n)=f(n-1)+2n-1$ for perfect squares
        \end{itemize}
\end{itemize}

\subsection*{Types of Recursive Definitions}
\begin{itemize}[leftmargin=*, itemsep=0pt]
    \item \textbf{First Order:}
    \begin{itemize}[leftmargin=*, itemsep=0pt]
        \item Base case: $f(1)$ given
        \item Recursive step: $f(n)$ in terms of $f(n-1)$
    \end{itemize}
    \item \textbf{k-th Order:}
    \begin{itemize}[leftmargin=*, itemsep=0pt]
        \item Base cases: $f(1), f(2), ..., f(k)$ given
        \item Uses $k$ previous terms
        \item $f(n)$ depends on $f(n-1), ..., f(n-k)$
    \end{itemize}
\end{itemize}

\subsection*{Famous Sequences}
\textbf{Fibonacci Sequence:}
\begin{itemize}[leftmargin=*, itemsep=0pt]
    \item Base cases: $f(1)=1$, $f(2)=1$
    \item Recursive: $f(n) = f(n-1) + f(n-2)$
    \item First terms: 1, 1, 2, 3, 5, 8, 13, 21, ...
    \item Explicit: $f(n) = \frac{1}{\sqrt{5}}[(1+\sqrt{5})^n - (1-\sqrt{5})^n]$
\end{itemize}

\textbf{Factorial:}
\begin{itemize}[leftmargin=*, itemsep=0pt]
    \item Base case: $f(1)=1$
    \item Recursive: $f(n) = n \cdot f(n-1)$
    \item Explicit: $f(n) = n!$
    \item First terms: 1, 2, 6, 24, 120, ...
\end{itemize}

\subsection*{Finding Explicit Forms}
\textbf{Inductive Reasoning Method:}
\begin{itemize}[leftmargin=*, itemsep=0pt]
    \item Calculate first several terms
    \item Look for patterns
    \item Form hypothesis for explicit form
    \item Prove using mathematical induction
\end{itemize}

\subsection*{Proof Methods}
\textbf{Mathematical Induction:}
\begin{itemize}[leftmargin=*, itemsep=0pt]
    \item Prove base case(s)
    \item Assume true for $k$
    \item Prove true for $k+1$
    \item Conclude true for all $n$
\end{itemize}

\textbf{Direct Substitution:}
\begin{itemize}[leftmargin=*, itemsep=0pt]
    \item Verify base case(s)
    \item Substitute explicit form into recursive definition
    \item Show equality holds
\end{itemize}

\subsection*{Recursive Trees}
\begin{itemize}[leftmargin=*, itemsep=0pt]
    \item Visual representation of recursive calls
    \item Shows dependency structure
    \item Useful for:
    \begin{itemize}[leftmargin=*, itemsep=0pt]
        \item Understanding computation flow
        \item Identifying repeated calculations
        \item Analyzing efficiency
    \end{itemize}
\end{itemize}

\end{multicols*}
\end{document}
